%=================================================================
% SAD - Documento de Arquitectura de Software
% Proyecto Tickets-NFT B2B (Stellar/Soroban)
%=================================================================
\documentclass[12pt,oneside]{article}

% ====== Paquetes y formato ======
\usepackage[spanish,es-nodecimaldot]{babel}
\usepackage[utf8]{inputenc}
\usepackage[T1]{fontenc}
\usepackage[a4paper,margin=2.5cm]{geometry}
\usepackage{setspace}
\onehalfspacing
\usepackage{graphicx}
\usepackage{float}
\usepackage{hyperref}
\hypersetup{
  colorlinks=true,
  linkcolor=blue,
  urlcolor=blue,
  citecolor=black,
  pdfauthor={Grupo 12},
  pdftitle={SAD - Proyecto Tickets-NFT B2B}
}
\usepackage{booktabs}
\usepackage{tabularx}
\usepackage{longtable}
\usepackage{array}
\usepackage{enumitem}
\setlist[itemize]{left=0pt,label=--,itemsep=2pt,topsep=4pt}
\setlist[enumerate]{left=0pt,itemsep=2pt,topsep=4pt}
\setcounter{secnumdepth}{3}
\setcounter{tocdepth}{3}
\usepackage{placeins}   % \FloatBarrier — evita que floats crucen secciones

% Evita páginas casi vacías por estiramiento vertical
\raggedbottom
% Penaliza líneas huérfanas/viudas (una sola línea al inicio/fin de página)
\widowpenalty=10000
\clubpenalty=10000
\usepackage{listings}
\usepackage{xcolor}
\usepackage{amssymb}

% Configuración de listings para código
\lstset{
  basicstyle=\ttfamily\small,
  breaklines=true,
  frame=single,
  numbers=left,
  numberstyle=\tiny\color{gray},
  keywordstyle=\color{blue},
  commentstyle=\color{green!60!black},
  stringstyle=\color{red!70!black},
  tabsize=2
}

% ====== Metadatos del proyecto ======
\newcommand{\projectname}{Reventa B2B de tickets con NFTs en Stellar (Soroban)}
\newcommand{\shortname}{Proyecto Tickets-NFT B2B}
\newcommand{\version}{0.1}
\newcommand{\docdate}{\today}
\newcommand{\course}{Trabajo de Grado -- Ingeniería de Sistemas}
\newcommand{\team}{Grupo 12}

% ====== Documento ======
\begin{document}

% ====== Portada ======
\begin{titlepage}
  \centering
  {\Large Documento de Especificación del Diseño\par}
  \vspace{1.2cm}
  {\huge \textbf{\projectname}\par}
  \vspace{0.3cm}
  {\Large (\shortname)\par}
  \vspace{1.2cm}
  {\large Versión 0.6 — \docdate\par}
  \vspace{0.8cm}
  \includegraphics[width=4cm]{pontificia-universidad-javeriana-seeklogo.png}\par
  \vspace{0.4cm}
  {\Large Pontificia Universidad Javeriana Bogotá\par}
  \vfill
  {\large \course\par}
  \vspace{0.3cm}
  {\large Director: Rafael Vicente Paéz Méndez, PhD\par}
  \vspace{0.3cm}
  {\large Cliente: B.A.F: Blockchain Acceleration Foundation\par}
  \vfill
  {\large\textbf{Autores:}\par}
  {\large
    Juan Diego Arias Durán\\
    Santiago Andrés Mesa Niño\\
    Andrés Felipe Manosalva Garzón\\
    Juan Camilo Carvajal Rodríguez\par}
  \vspace{1cm}
\end{titlepage}

% ====== Preliminares ======
\pagenumbering{roman}

% ------ Historial de cambios ------
\section*{Historial de cambios}
\begin{longtable}{@{}p{2cm} p{3cm} p{7.5cm} p{2.5cm}@{}}
  \toprule
  \textbf{Versión} & \textbf{Fecha} & \textbf{Descripción} & \textbf{Autor(es)}\\
  \midrule
  \endfirsthead
  \toprule
  \textbf{Versión} & \textbf{Fecha} & \textbf{Descripción} & \textbf{Autor(es)}\\
  \midrule
  \endhead
  \bottomrule
  \endfoot
  0.1 & 10/03/2026 & Versión inicial del SAD: estructura base del documento. & \team\\
  0.2 & 13/03/2026 & Desarrollo de las secciones 1 (Introducción), 2 (Representación arquitectónica) y 3 (Objetivos y restricciones arquitectónicas). & \team\\
  0.3 & 15/03/2026 & Desarrollo de las secciones 4 (Vista de casos de uso), 5 (Vista lógica) y 6 (Vista de procesos). & \team\\
  0.4 & 24/03/2026 & Actualización integral: alineación con la implementación real de los contratos (nombres en español, versionado burn/remint, sistema de verificadores, invalidación on-chain), corrección de la arquitectura Frontend$\rightarrow$Backend$\rightarrow$Blockchain, eliminación de App de Escaneo y modo offline no implementados, actualización de diagramas C4. & \team\\
  0.5 & 25/04/2026 & Sincronización con estado final del prototipo: corrección de stack backend (Express en lugar de NestJS), documentación de integración Soroban completa (venta primaria, reventa atómica, asegurar on-chain, mercado secundario P2P, cancelación de listado), incorporación de Freighter Wallet como componente implementado, actualización de modelo de datos off-chain con campos Web3, adición de rol STAFF y panel de administración (Phase 4), corrección de flujo de verificación (escáner QR implementado, sin modo offline), actualización de rutas frontend y variables de entorno. & \team\\
  0.6 & 11/05/2026 & Referencias bibliográficas movidas de la sección 1.4 al final del documento como sección independiente. Diagrama de casos de uso actualizado: solo Admin Plataforma puede crear eventos y desplegar contratos inteligentes. Actor renombrado de Comprador a Usuario en todos los casos de uso. & \team\\
\end{longtable}

\newpage

% ------ Índices ------
\tableofcontents
\newpage
\listoffigures
\listoftables
\clearpage

% ====== Cuerpo ======
\pagenumbering{arabic}

%===============================================================
\section{Introducción}
\label{sec:introduccion}
%===============================================================

\subsection{Propósito}
\label{subsec:proposito}

Este documento describe la \textbf{arquitectura de software} del sistema \textbf{\shortname}. Su propósito es:

\begin{itemize}
  \item Comunicar las decisiones arquitectónicas fundamentales del sistema a todos los interesados.
  \item Servir como referencia técnica para el equipo de desarrollo durante la implementación.
  \item Documentar las vistas arquitectónicas que permiten comprender la estructura, comportamiento y despliegue del sistema.
  \item Justificar las decisiones de diseño y los \textit{trade-offs} realizados.
\end{itemize}

\subsection{Alcance}
\label{subsec:alcance}

El SAD abarca la arquitectura del prototipo académico del sistema de reventa B2B de tickets con NFTs sobre Stellar (Soroban), incluyendo:

\begin{itemize}
  \item El contrato inteligente Soroban (lógica on-chain).
  \item El backend/API que interactúa con el contrato.
  \item El frontend web de demostración.
  \item Las integraciones con la red Stellar (testnet).
\end{itemize}

\subsection{Definiciones, acrónimos y abreviaturas}
\label{subsec:definiciones}

\begin{longtable}{@{}p{3.5cm} p{11cm}@{}}
  \toprule
  \textbf{Término} & \textbf{Definición}\\
  \midrule
  \endhead
  \bottomrule
  \endfoot
  SAD & \textit{Software Architecture Document}. Este documento.\\
  NFT & \textit{Non-Fungible Token}. Activo digital único que representa un ticket en blockchain.\\
  Stellar & Red blockchain orientada a pagos y emisión de activos digitales.\\
  Soroban & Plataforma de contratos inteligentes sobre Stellar, basada en WebAssembly (Wasm).\\
  WASM & \textit{WebAssembly}. Formato de bytecode al que se compilan los contratos Soroban.\\
  B2B & \textit{Business to Business}. Modelo de negocio del sistema.\\
  API & \textit{Application Programming Interface}.\\
  REST & \textit{Representational State Transfer}. Estilo arquitectónico para APIs web.\\
  Testnet & Red de pruebas de Stellar sin valor monetario real.\\
  SRS & \textit{Software Requirements Specification}. Documento de requerimientos del proyecto.\\
  SPMP & \textit{Software Project Management Plan}. Plan de gestión del proyecto.\\
\end{longtable}

\subsection{Visión general del documento}
\label{subsec:vision-general}

El resto del documento se organiza así:

\begin{itemize}
  \item \textbf{Sección~\ref{sec:representacion}}: describe el modelo arquitectónico utilizado y las vistas que se presentan.
  \item \textbf{Sección~\ref{sec:objetivos-restricciones}}: enumera los objetivos y restricciones arquitectónicas.
  \item \textbf{Sección~\ref{sec:vista-casos-uso}}: vista de casos de uso.
  \item \textbf{Sección~\ref{sec:vista-logica}}: vista lógica (paquetes, clases, capas).
  \item \textbf{Sección~\ref{sec:vista-procesos}}: vista de procesos (concurrencia, flujos).
  \item \textbf{Sección~\ref{sec:vista-despliegue}}: vista de despliegue (infraestructura física/virtual).
  \item \textbf{Sección~\ref{sec:vista-implementacion}}: vista de implementación (componentes, módulos).
  \item \textbf{Sección~\ref{sec:vista-datos}}: vista de datos (modelo de datos on-chain y off-chain).
  \item \textbf{Sección~\ref{sec:calidad}}: atributos de calidad y tácticas arquitectónicas.
  \item \textbf{Sección~\ref{sec:decisiones}}: registro de decisiones arquitectónicas (ADRs).
\end{itemize}

%===============================================================
\FloatBarrier
\clearpage
\section{Representación arquitectónica}
\label{sec:representacion}
%===============================================================

La arquitectura de Stellar Tickets se describe mediante el modelo \textbf{C4} de Simon Brown, que provee niveles de abstracción progresivos (contexto, contenedores, componentes, código) para comunicar la arquitectura a audiencias con distintos niveles de detalle técnico. Este modelo es especialmente adecuado para sistemas híbridos como el presente, donde la lógica se distribuye entre contratos inteligentes on-chain y servicios convencionales off-chain.

\subsection{Niveles C4 utilizados}

Se utilizan los dos primeros niveles del modelo C4:

\subsubsection{Nivel 1 -- Contexto del sistema}

En el modelo C4, el diagrama de contexto es el nivel de abstracción más alto. Su objetivo es responder a la pregunta: \textit{¿qué es el sistema y quién interactúa con él?} No muestra detalles técnicos internos, sino que ubica al sistema como una caja negra dentro de su entorno, identificando las personas que lo usan y los sistemas externos de los que depende o con los que se comunica.

La Figura~\ref{fig:c4-contexto} presenta el diagrama de contexto de Stellar Tickets.

\begin{figure}[htbp]
  \centering
  \includegraphics[width=0.7\textwidth]{contextoSecureTicket.png}
  \caption{Diagrama C4 Nivel 1 -- Contexto del sistema Stellar Tickets}
  \label{fig:c4-contexto}
\end{figure}

\paragraph{Interpretación del diagrama.} El sistema central es la \textbf{plataforma Stellar Tickets}, compuesta por un frontend web (React SPA), un backend API (Express) con su Indexador Soroban, y los contratos inteligentes Soroban. Alrededor se identifican los siguientes actores y sistemas externos:

\begin{itemize}
  \item \textbf{Organizador} (empresa de boletería): configura eventos, define comisiones y gestiona boletos a través del panel de administración web. Sus operaciones on-chain son orquestadas por el backend.
  \item \textbf{Comprador} (usuario final): navega eventos, compra boletos en venta primaria o reventa, y gestiona su inventario de boletos a través de la aplicación web TuTicket.
  \item \textbf{Verificador} (personal del evento): persona autorizada por el organizador para validar boletos en la puerta del evento. Utiliza la funcionalidad de verificación del sistema para marcar boletos como redimidos on-chain.
  \item \textbf{Administrador de la plataforma} (equipo técnico): despliega los contratos en la red Stellar mediante el Factory, configura el WASM hash y mantiene la infraestructura operativa.
  \item \textbf{Blockchain Stellar}: red pública descentralizada que almacena el estado de los contratos de forma inmutable. Los contratos se ejecutan sobre Soroban (motor de smart contracts basado en WASM).
  \item \textbf{SAC XLM Nativo}: Stellar Asset Contract del token XLM nativo de la red Stellar. El contrato de evento lo invoca para ejecutar transferencias atómicas de pago durante compras y reventas.
  \item \textbf{Freighter Wallet}: extensión de navegador que custodia la clave privada del usuario. El frontend la usa para solicitar la firma de transacciones Soroban sin exponer la clave al servidor.
\end{itemize}

El límite del sistema define qué está bajo control del equipo de desarrollo (frontend, backend, contratos) y qué es externo (red Stellar, SAC XLM, Freighter).

\subsubsection{Nivel 2 -- Contenedores}

El diagrama de contenedores es el segundo nivel del modelo C4. Responde a la pregunta: \textit{¿cuáles son las piezas técnicas que componen el sistema y cómo se comunican entre sí?} Un ``contenedor'' en C4 no se refiere a Docker, sino a una unidad desplegable independientemente: una aplicación web, una API, una base de datos, un contrato inteligente, etc.

La Figura~\ref{fig:c4-contenedores} presenta el diagrama de contenedores de Stellar Tickets.

\begin{figure}[htbp]
  \centering
  \includegraphics[width=0.6\textwidth]{02_c4_contenedores.png}
  \caption{Diagrama C4 Nivel 2 -- Contenedores del sistema Stellar Tickets}
  \label{fig:c4-contenedores}
\end{figure}

\paragraph{Interpretación del diagrama.} El sistema se descompone en los siguientes contenedores, organizados en cuatro zonas:

\textbf{Zona Cliente} (lo que ejecuta el usuario en su dispositivo):
\begin{itemize}
  \item \textbf{Aplicación Web TuTicket}: aplicación frontend (React + Vite + TypeScript) que permite a compradores navegar eventos, comprar boletos, gestionar su inventario y realizar reventas. Incluye un panel de administración para organizadores. Se comunica exclusivamente con el Backend API mediante llamadas REST.
\end{itemize}

\textbf{Zona Servidor} (lo que se ejecuta en la infraestructura de la plataforma):
\begin{itemize}
  \item \textbf{Backend API} (\texttt{Express 4 + TypeScript}): servidor Express monolítico (\texttt{src/server.ts}) que implementa la lógica de negocio off-chain: autenticación JWT (bcrypt + jsonwebtoken), gestión de eventos, carrito de compras, checkout fiat, emisión de boletos y endpoints de administración. Actúa como \textit{proxy XDR}: construye transacciones Soroban sin firmar y las retorna al frontend para que Freighter las firme. Firma server-side las operaciones del organizador (\texttt{crear\_boleto}, \texttt{listar\_boleto}, \texttt{cancelar\_venta}).
  \item \textbf{Indexador Soroban} (\texttt{src/indexer.ts}): proceso en segundo plano que hace polling al RPC de Soroban cada 5 segundos. Sincroniza 7 tipos de eventos on-chain hacia PostgreSQL. Solo corre en entornos long-lived (no en Vercel).
  \item \textbf{PostgreSQL (Supabase)}: base de datos relacional que persiste el estado híbrido Web2+Web3: usuarios, eventos, boletos con campos on-chain (\texttt{contract\_address}, \texttt{ticket\_root\_id}, \texttt{version}, \texttt{is\_for\_sale}, \texttt{resale\_price}, \texttt{owner\_wallet}), órdenes y pagos.
  \item \textbf{Nodo RPC Soroban}: punto de acceso a la red Stellar (\texttt{soroban-testnet.stellar.org}). Permite al backend enviar transacciones firmadas y al Indexador consultar eventos on-chain.
\end{itemize}

\textbf{Zona Blockchain} (Stellar Testnet, lo que se ejecuta de forma descentralizada):
\begin{itemize}
  \item \textbf{Contrato Factory} (\texttt{FabricaBoletos}): contrato Soroban que despliega un contrato de evento independiente por cada evento. Garantiza aislamiento de configuración y fondos entre eventos.
  \item \textbf{Contrato de Evento} (\texttt{ContratoEvento}, 11 instancias desplegadas): contrato Soroban desplegado por evento que contiene toda la lógica de negocio on-chain: creación de boletos, venta primaria, reventa atómica con burn/remint, distribución de comisiones y gestión de verificadores. Es el componente más crítico del sistema.
  \item \textbf{SAC XLM Nativo}: Stellar Asset Contract del token XLM nativo. El contrato de evento lo invoca para ejecutar transferencias atómicas de pago.
\end{itemize}

\textbf{Comunicaciones clave}:
\begin{itemize}
  \item La Aplicación Web TuTicket se comunica con el Backend API mediante REST/JSON para todas las operaciones (autenticación, navegación, compra, gestión de boletos).
  \item El Backend API interactúa con PostgreSQL para persistencia off-chain y con el Nodo RPC Soroban para operaciones on-chain.
  \item El Contrato de Evento invoca al SAC XLM nativo para procesar transferencias de pago de forma atómica dentro de la misma transacción Soroban.
  \item La Red Stellar subyace como infraestructura que persiste el estado de todos los contratos de forma inmutable y distribuida.
\end{itemize}

\textbf{Estado de integración:} la conexión Backend API $\leftrightarrow$ Soroban está completamente implementada. El backend (Express + TypeScript) ejecuta operaciones on-chain reales sobre la Stellar Testnet: creación de boletos (\texttt{crear\_boleto}), listado para reventa (\texttt{listar\_boleto}), cancelación (\texttt{cancelar\_venta}), construcción de transacciones XDR para que el comprador las firme con Freighter (\texttt{comprar\_boleto}), y envío de transacciones firmadas. El Indexador sincroniza eventos Soroban hacia PostgreSQL cada 5 segundos.

\subsection{Mapeo de vistas a secciones del documento}

Las secciones restantes del SAD profundizan en aspectos específicos de la arquitectura, complementando los diagramas C4:

\begin{table}[htbp]
  \centering
  \begin{tabularx}{\textwidth}{@{}l l X@{}}
    \toprule
    \textbf{Vista} & \textbf{Sección} & \textbf{Contenido principal}\\
    \midrule
    Casos de uso & \ref{sec:vista-casos-uso} & Casos de uso arquitectónicamente significativos y diagrama C4 de contexto.\\
    Lógica & \ref{sec:vista-logica} & Descomposición en capas y módulos; diagrama C4 de contenedores; separación on-chain/off-chain.\\
    Procesos & \ref{sec:vista-procesos} & Flujos de interacción principales: compra primaria, reventa atómica, redención y verificación.\\
    Despliegue & \ref{sec:vista-despliegue} & Topología de nodos: presentación, servicios, datos y red Stellar testnet.\\
    Implementación & \ref{sec:vista-implementacion} & Organización del código fuente, artefactos compilados (WASM) y dependencias.\\
    Datos & \ref{sec:vista-datos} & Modelo de almacenamiento on-chain (Soroban) y off-chain (PostgreSQL).\\
    \bottomrule
  \end{tabularx}
  \caption{Mapeo entre vistas arquitectónicas y secciones del SAD}
  \label{tab:mapeo-vistas}
\end{table}

Adicionalmente, el documento incluye una sección de \textbf{atributos de calidad} (Sección~\ref{sec:calidad}) con escenarios y tácticas arquitectónicas, y un \textbf{registro de decisiones arquitectónicas} (Sección~\ref{sec:decisiones}) en formato ADR.

%===============================================================
\FloatBarrier
\clearpage
\section{Objetivos y restricciones arquitectónicas}
\label{sec:objetivos-restricciones}
%===============================================================

Esta sección presenta los \textit{drivers} arquitectónicos: los atributos de calidad prioritarios y las restricciones que condicionan las decisiones de diseño del sistema.

\subsection{Objetivos arquitectónicos}
\label{subsec:objetivos}

Los siguientes atributos de calidad, derivados del SRS y de las necesidades del modelo B2B, guían las decisiones arquitectónicas:

\begin{enumerate}
  \item \textbf{Trazabilidad e inmutabilidad.} Cada ticket debe tener un historial de propiedad verificable e inalterable. Esto motiva el uso de blockchain como fuente de verdad para la propiedad y las transferencias, complementado con un indexador off-chain para consultas rápidas.

  \item \textbf{Atomicidad en transacciones de reventa.} La operación de compra en el mercado secundario debe ser atómica: la transferencia del ticket y la distribución de pagos (vendedor, comisión organizador, comisión plataforma) ocurren en una sola transacción o no ocurren en absoluto. Esto elimina estados inconsistentes donde el pago se procesa pero el ticket no se transfiere, o viceversa.

  \item \textbf{Seguridad y autenticidad.} La arquitectura debe garantizar que los tickets no puedan falsificarse ni duplicarse. La custodia de claves es no-custodial (Freighter Wallet), de modo que solo el propietario legítimo puede autorizar transferencias mediante firma criptográfica.

  \item \textbf{Adaptabilidad B2B (multi-ticketera).} El sistema debe permitir que múltiples ticketeras operen de forma independiente, cada una con sus propios eventos, políticas de comisión y configuración. Esto motiva la arquitectura de \textit{contrato por evento} desplegado por un Factory. El frontend TuTicket simula una ticketera cliente de la plataforma.

  \item \textbf{Disponibilidad de consultas.} Aunque la blockchain garantiza persistencia, su modelo de lectura es limitado. El sistema debe ofrecer consultas rápidas sobre estado de tickets, historial y auditoría sin depender exclusivamente de llamadas on-chain. Esto justifica la capa de indexación y caché off-chain respaldada por PostgreSQL.

  \item \textbf{Verificabilidad en puerta.} La validación de boletos en el punto de acceso al evento se realiza mediante verificadores autorizados por el organizador. El contrato Soroban implementa un registro de verificadores (\texttt{Verificador(Address)}) que restringe quién puede marcar un boleto como redimido, garantizando que la redención represente un hecho físico de asistencia.
\end{enumerate}

\subsection{Restricciones arquitectónicas}
\label{subsec:restricciones}

\subsubsection{Restricciones tecnológicas}

\begin{itemize}
  \item \textbf{Stellar y Soroban:} la lógica on-chain debe implementarse como contratos inteligentes Soroban, compilados a WebAssembly (WASM) y escritos en Rust. Esto impone el uso del \texttt{soroban-sdk} y sus limitaciones (modelo de storage basado en claves, sin eventos complejos nativos, tamaño limitado de transacciones).
  \item \textbf{Testnet:} el prototipo opera exclusivamente sobre la testnet de Stellar, lo que implica que los activos no tienen valor monetario real y que la red puede presentar interrupciones o reinicios periódicos.
  \item \textbf{Token de pago:} las transacciones de compra y reventa utilizan \textbf{XLM nativo} (Lumens) como token de pago. Los valores se expresan en \textit{stroops} (1 XLM = 10\textsuperscript{7} stroops). El Stellar Asset Contract (SAC) nativo (\texttt{CDLZFC3S...}) es el contrato de token que invoca el \texttt{ContratoEvento}.
  \item \textbf{Firma de transacciones:} operaciones iniciadas por el organizador (\texttt{crear\_boleto}, \texttt{listar\_boleto}, \texttt{cancelar\_venta}) son firmadas server-side con la clave del organizador. La operación de compra (\texttt{comprar\_boleto}) es firmada por el comprador a través de \textbf{Freighter Wallet} (extensión de navegador no-custodial): el backend construye la transacción XDR sin firmar, el frontend la pasa a Freighter para firma, y luego el backend la envía a la red Stellar.
\end{itemize}

\subsubsection{Restricciones organizacionales}

\begin{itemize}
  \item \textbf{Equipo reducido:} 4 desarrolladores con dedicación parcial (proyecto académico), lo que favorece una arquitectura con separación clara de responsabilidades que permita trabajo en paralelo (contrato, backend, frontend, verificación).
  \item \textbf{Tiempo acotado:} dos semestres académicos ($\sim$32 semanas), lo que limita el alcance a un prototipo funcional demostrable.
  \item \textbf{Contexto académico:} el sistema debe ser documentable, demostrable y evaluable por docentes, priorizando claridad arquitectónica sobre optimización de producción.
\end{itemize}

\subsubsection{Restricciones de integración}

\begin{itemize}
  \item \textbf{Modelo B2B no invasivo:} el sistema se ofrece como complemento a plataformas de boletería existentes, no como reemplazo. La integración se realiza vía API REST, sin requerir modificaciones en los sistemas internos de la ticketera.
  \item \textbf{Separación on-chain/off-chain:} las operaciones de escritura que afectan propiedad y pagos ocurren on-chain (fuente de verdad), mientras que las consultas de lectura, auditoría e indexación se resuelven off-chain para rendimiento y usabilidad.
\end{itemize}

%===============================================================
\FloatBarrier
\clearpage
\section{Vista de casos de uso}
\label{sec:vista-casos-uso}
%===============================================================

Esta sección identifica los casos de uso arquitectónicamente significativos del sistema: aquellos que ejercitan los principales componentes, imponen restricciones de diseño o exponen los atributos de calidad prioritarios descritos en la Sección~\ref{sec:objetivos-restricciones}.

\subsection{Diagrama de casos de uso}

La Figura~\ref{fig:casos-uso} presenta los casos de uso principales del sistema y sus relaciones con los actores identificados en el diagrama C4 de contexto (Figura~\ref{fig:c4-contexto}).

\begin{figure}[htbp]
  \centering
  \includegraphics[width=0.7\textwidth]{CasosDeUso_SAD.png}
  \caption{Diagrama de casos de uso arquitectónicamente significativos}
  \label{fig:casos-uso}
\end{figure}

\subsection{Descripción de casos de uso}

A continuación se describe cada caso de uso, destacando su relevancia arquitectónica.

\clearpage
\subsubsection{CU-01: Crear evento y desplegar contrato}

\begin{table}[h!]
  \centering
  \begin{tabularx}{\textwidth}{@{}p{3.5cm} X@{}}
    \toprule
    \textbf{Actores} & Admin Plataforma\\
    \midrule
    \textbf{Descripción} & El administrador de la plataforma configura un nuevo evento (nombre, fecha, wallets de comisión, políticas de reventa). La plataforma invoca el Factory para desplegar un contrato Soroban específico para ese evento en la testnet.\\
    \midrule
    \textbf{Precondiciones} & El administrador está autenticado y tiene permisos de creación de eventos.\\
    \midrule
    \textbf{Flujo principal} &
      1. El administrador ingresa los datos del evento y configura fees.\newline
      2. El backend valida los parámetros y registra el evento en PostgreSQL.\newline
      3. El backend invoca el contrato Factory en Stellar.\newline
      4. El Factory despliega un nuevo contrato de evento con los parámetros configurados.\newline
      5. El sistema registra la dirección del contrato desplegado.\\
    \midrule
    \textbf{Postcondiciones} & Existe un contrato de evento independiente en la testnet, listo para emitir tickets.\\
    \midrule
    \textbf{Relevancia arquitectónica} & Ejerce el patrón Factory, la separación on-chain/off-chain y la adaptabilidad multi-ticketera. Cada evento opera con un contrato aislado, lo que garantiza independencia de configuración y fondos.\\
    \bottomrule
  \end{tabularx}
  \caption{CU-01: Crear evento y desplegar contrato}
  \label{tab:cu01}
\end{table}

\clearpage
\subsubsection{CU-02: Compra primaria de ticket}

\begin{table}[h!]
  \centering
  \begin{tabularx}{\textwidth}{@{}p{3.5cm} X@{}}
    \toprule
    \textbf{Actores} & Usuario\\
    \midrule
    \textbf{Descripción} & El usuario adquiere un ticket emitido por el organizador. El 100\% del pago se transfiere al organizador.\\
    \midrule
    \textbf{Precondiciones} & El evento existe, tiene tickets disponibles y el usuario tiene fondos suficientes en su wallet.\\
    \midrule
    \textbf{Flujo principal} &
      1. El usuario selecciona un boleto disponible en la aplicación web TuTicket.\newline
      2. El frontend envía la solicitud de compra al Backend API.\newline
      3. El backend orquesta la invocación de \texttt{comprar\_boleto} en el contrato de evento.\newline
      4. El contrato transfiere el 100\% del pago al organizador mediante \texttt{ClienteToken} y asigna el boleto al usuario, marcando \texttt{es\_reventa = true}.\newline
      5. El backend registra la operación en PostgreSQL.\\
    \midrule
    \textbf{Postcondiciones} & El usuario es el nuevo propietario del boleto. El pago fue transferido al organizador. Futuras ventas de este boleto serán tratadas como reventa (con comisiones y burn/remint).\\
    \midrule
    \textbf{Relevancia arquitectónica} & Ejerce la orquestación Backend$\rightarrow$Soroban, la escritura on-chain con registro off-chain posterior, y la transición \texttt{es\_reventa = false $\rightarrow$ true} que determina el comportamiento de todas las transacciones futuras sobre ese boleto.\\
    \bottomrule
  \end{tabularx}
  \caption{CU-02: Compra primaria de ticket}
  \label{tab:cu02}
\end{table}

\clearpage
\subsubsection{CU-03: Listar ticket para reventa}

\begin{table}[h!]
  \centering
  \begin{tabularx}{\textwidth}{@{}p{3.5cm} X@{}}
    \toprule
    \textbf{Actores} & Usuario\\
    \midrule
    \textbf{Descripción} & El propietario de un ticket lo publica en el marketplace para reventa, estableciendo un precio.\\
    \midrule
    \textbf{Precondiciones} & El usuario es propietario del ticket, el ticket está vigente (no consumido ni invalidado) y el evento permite reventa.\\
    \midrule
    \textbf{Flujo principal} &
      1. El usuario selecciona un boleto de su inventario y fija precio de reventa.\newline
      2. El frontend envía la solicitud al Backend API.\newline
      3. El backend invoca \texttt{listar\_boleto} en el contrato de evento, que marca \texttt{en\_venta = true} y actualiza el precio.\newline
      4. El marketplace muestra el boleto disponible para compra.\\
    \midrule
    \textbf{Postcondiciones} & El boleto está visible en el marketplace con estado \texttt{en\_venta = true}.\\
    \midrule
    \textbf{Relevancia arquitectónica} & Ejerce la transición de estados del ticket on-chain y la sincronización con el caché off-chain para consultas del marketplace.\\
    \bottomrule
  \end{tabularx}
  \caption{CU-03: Listar ticket para reventa}
  \label{tab:cu03}
\end{table}

\clearpage
\subsubsection{CU-04: Compra en reventa atómica}

\begin{table}[h!]
  \centering
  \begin{tabularx}{\textwidth}{@{}p{3.5cm} X@{}}
    \toprule
    \textbf{Actores} & Usuario\\
    \midrule
    \textbf{Descripción} & El usuario adquiere un ticket en el mercado secundario. En una sola transacción atómica se ejecutan: pago al vendedor, comisión al organizador, comisión a la plataforma, burn de la versión anterior y mint de una nueva versión del ticket.\\
    \midrule
    \textbf{Precondiciones} & El ticket está listado para reventa, el usuario tiene fondos suficientes y no es el propietario actual.\\
    \midrule
    \textbf{Flujo principal} &
      1. El usuario selecciona un boleto en reventa y confirma la compra.\newline
      2. El frontend envía la solicitud al Backend API.\newline
      3. El backend invoca \texttt{comprar\_boleto} en el contrato de evento, que ejecuta atómicamente:\newline
      \hspace{1em}a) Transferencia de comisión al organizador (\texttt{ComisionOrganizador\%}).\newline
      \hspace{1em}b) Transferencia de comisión a la plataforma (\texttt{ComisionPlataforma\%}).\newline
      \hspace{1em}c) Transferencia del pago neto al vendedor.\newline
      \hspace{1em}d) Burn: la versión actual se marca como \texttt{invalidado = true}.\newline
      \hspace{1em}e) Remint: se crea nueva versión (\texttt{version + 1}) asignada al usuario.\newline
      4. El contrato emite evento \texttt{BoletoRevendido} on-chain indexado por \texttt{ticket\_root\_id}.\newline
      5. El backend actualiza PostgreSQL.\\
    \midrule
    \textbf{Flujo alternativo} & Si cualquier paso falla, toda la transacción se revierte. No existe estado intermedio.\\
    \midrule
    \textbf{Postcondiciones} & El usuario posee la nueva versión del ticket. Los pagos y comisiones fueron distribuidos. La versión anterior está invalidada.\\
    \midrule
    \textbf{Relevancia arquitectónica} & Es el caso de uso más crítico del sistema. Ejerce la atomicidad transaccional de Soroban, el patrón burn+remint, la distribución de comisiones embebida en el contrato y la trazabilidad por \texttt{ticket\_root\_id}. Valida el driver de atomicidad e inmutabilidad.\\
    \bottomrule
  \end{tabularx}
  \caption{CU-04: Compra en reventa atómica}
  \label{tab:cu04}
\end{table}

\clearpage
\subsubsection{CU-05: Verificar ticket en puerta}

\begin{table}[h!]
  \centering
  \begin{tabularx}{\textwidth}{@{}p{3.5cm} X@{}}
    \toprule
    \textbf{Actores} & Personal del Evento\\
    \midrule
    \textbf{Descripción} & Un verificador autorizado por el organizador valida el boleto del asistente en la puerta del evento y lo marca como redimido de forma irreversible en el contrato.\\
    \midrule
    \textbf{Precondiciones} & El boleto existe, no está invalidado, no ha sido redimido previamente, y el verificador está registrado como autorizado en el contrato (\texttt{es\_verificador = true}).\\
    \midrule
    \textbf{Flujo principal} &
      1. El asistente presenta su boleto (código QR) en la puerta del evento. El QR contiene un payload JSON con el \texttt{ticketId} (UUID en PostgreSQL).\newline
      2. El verificador (rol \texttt{STAFF} o \texttt{ADMIN}) abre la página \texttt{/escanear} en su dispositivo.\newline
      3. La cámara del dispositivo lee el QR. El frontend envía \texttt{POST /api/admin/scan \{ticketId\}} al Backend API.\newline
      4. El backend valida que el boleto existe en PostgreSQL con estado \texttt{ACTIVE}.\newline
      5. Si es válido, actualiza el estado a \texttt{CANCELLED} en PostgreSQL y retorna \texttt{\{success: true\}}.\newline
      6. La interfaz muestra ``Acceso Permitido'' (verde) al verificador.\\
    \midrule
    \textbf{Postcondiciones} & El boleto queda marcado como \texttt{CANCELLED} en PostgreSQL. No puede volver a escanearse.\\
    \midrule
    \textbf{Relevancia arquitectónica} & Ejerce el sistema de roles off-chain (\texttt{STAFF}/\texttt{ADMIN}), la validación de estado en PostgreSQL y la decisión de diseño de scan DB-only. El contrato Soroban implementa \texttt{redimir\_boleto} y el sistema de verificadores on-chain como mecanismo complementario de trazabilidad criptográfica, activable en producción donde la latencia adicional sea aceptable.\\
    \bottomrule
  \end{tabularx}
  \caption{CU-05: Verificar ticket en puerta}
  \label{tab:cu05}
\end{table}

\clearpage
\subsubsection{CU-06: Invalidar boleto}

\begin{table}[h!]
  \centering
  \begin{tabularx}{\textwidth}{@{}p{3.5cm} X@{}}
    \toprule
    \textbf{Actores} & Organizador\\
    \midrule
    \textbf{Descripción} & El organizador invalida un boleto de forma administrativa (evento cancelado, error de emisión, etc.), marcándolo como \texttt{invalidado = true} on-chain.\\
    \midrule
    \textbf{Precondiciones} & El boleto existe y no ha sido invalidado previamente. El actor es el organizador del evento.\\
    \midrule
    \textbf{Flujo principal} &
      1. El organizador solicita la invalidación del boleto desde el panel de administración.\newline
      2. El backend invoca \texttt{invalidar\_boleto} en el contrato de evento.\newline
      3. El contrato marca el boleto como \texttt{invalidado = true}, \texttt{en\_venta = false} (irreversible).\newline
      4. El contrato emite evento \texttt{BoletoInvalidadoEvt}.\newline
      5. El backend actualiza el estado en PostgreSQL.\\
    \midrule
    \textbf{Postcondiciones} & El boleto no puede ser listado, comprado, revendido ni redimido. El estado es terminal.\\
    \midrule
    \textbf{Relevancia arquitectónica} & Ejerce la invalidación on-chain (a diferencia de versiones anteriores donde se manejaba solo off-chain), la irreversibilidad de estado y la autorización exclusiva del organizador para operaciones administrativas.\\
    \bottomrule
  \end{tabularx}
  \caption{CU-06: Invalidar boleto}
  \label{tab:cu06}
\end{table}

%===============================================================
\FloatBarrier
\clearpage
\section{Vista lógica}
\label{sec:vista-logica}
%===============================================================

Esta sección describe la descomposición lógica del sistema en capas y módulos, profundizando el diagrama C4 de contenedores (Figura~\ref{fig:c4-contenedores}). La separación fundamental es entre la lógica \textbf{on-chain} (ejecutada de forma descentralizada en la red Stellar) y la lógica \textbf{off-chain} (ejecutada en infraestructura controlada por la plataforma).

\subsection{Arquitectura en capas}

El sistema se organiza en cuatro capas lógicas, cada una con responsabilidades claramente delimitadas:

\begin{enumerate}
  \item \textbf{Capa de presentación} (Frontend): interfaz de usuario que permite a los distintos actores interactuar con el sistema.
  \item \textbf{Capa de servicios} (Backend/API): lógica de negocio off-chain, orquestación de operaciones, autenticación y exposición de endpoints.
  \item \textbf{Capa de datos} (PostgreSQL + Caché): persistencia off-chain para consultas rápidas, auditoría e indexación.
  \item \textbf{Capa blockchain} (Stellar/Soroban): lógica de negocio on-chain, fuente de verdad para propiedad de tickets y pagos.
\end{enumerate}

\begin{table}[htbp]
  \centering
  \begin{tabularx}{\textwidth}{@{}p{3cm} p{4cm} X@{}}
    \toprule
    \textbf{Capa} & \textbf{Contenedores} & \textbf{Responsabilidad}\\
    \midrule
    Presentación & Aplicación Web TuTicket (React/Vite) & Renderizar interfaces, capturar acciones del usuario, comunicarse con el Backend API.\\
    Servicios & Backend API (Express + Indexador Soroban) & Autenticar usuarios (JWT/bcrypt), orquestar operaciones de negocio, gestionar carrito, procesar checkout, exponer API REST. Firma server-side de \texttt{crear\_boleto}/\texttt{listar\_boleto}/\texttt{cancelar\_venta}. Construcción de XDR para firma por Freighter. Sincronización de eventos on-chain mediante indexador en segundo plano.\\
    Datos & PostgreSQL (schemas: public, ticketing) & Persistir usuarios, eventos, sedes, boletos, órdenes, pagos y auditoría. Fuente de datos para consultas del frontend.\\
    Blockchain & Contrato Factory (\texttt{FabricaBoletos}), Contrato por Evento (\texttt{ContratoEvento}), SAC XLM nativo & Crear boletos, transferir propiedad, distribuir comisiones, redimir y invalidar boletos. Fuente de verdad inmutable. 11 contratos desplegados en Stellar Testnet.\\
    \bottomrule
  \end{tabularx}
  \caption{Descomposición en capas lógicas}
  \label{tab:capas}
\end{table}

\subsection{Capa de presentación}

La capa de presentación consiste en una aplicación web:

\subsubsection{Aplicación Web TuTicket (React + Vite + TypeScript)}

Aplicación SPA (\textit{Single Page Application}) que simula una ticketera cliente de la plataforma. Utiliza React 18 con Vite como bundler, TypeScript, Tailwind CSS para estilos y shadcn/ui (basado en Radix UI) como librería de componentes. Expone las siguientes funcionalidades:

\begin{itemize}
  \item \textbf{Comprador}: navegación de eventos por categoría, búsqueda, selección de boletos y asientos, carrito de compras, checkout en 3 pasos (datos del comprador $\rightarrow$ método de pago $\rightarrow$ confirmación), historial de compras y gestión de boletos propios.
  \item \textbf{Administrador}: panel de administración para gestión de eventos, creación de tipos de boletos, consulta de estadísticas de ventas y ocupación.
\end{itemize}

La comunicación con la blockchain \textbf{no es directa}: el frontend se comunica con el Backend API mediante llamadas REST/JSON autenticadas con JWT para la mayoría de operaciones. La excepción es la \textbf{firma de transacciones de compra en reventa}: el frontend recibe el XDR sin firmar del backend y lo pasa a la extensión \textbf{Freighter Wallet} del navegador para obtener la firma del comprador, sin que el backend tenga acceso a claves privadas del usuario.

\textbf{Rutas principales (23 rutas)}: \texttt{/} (inicio), \texttt{/eventos} (catálogo), \texttt{/eventos/:categoria}, \texttt{/buscar}, \texttt{/evento/:id} (detalle + marketplace P2P), \texttt{/evento/:id/boletas}, \texttt{/evento/:id/asientos}, \texttt{/carrito}, \texttt{/checkout}, \texttt{/confirmacion}, \texttt{/login}, \texttt{/registro}, \texttt{/mi-cuenta} (perfil, entradas, compras, \textbf{ventas P2P}), \texttt{/admin} (panel de administración, solo rol ADMIN), \texttt{/escanear} (escáner QR de verificación, roles ADMIN y STAFF).

\subsection{Capa de servicios (Backend)}

El backend es una aplicación \textbf{Express 4 + TypeScript} que centraliza la lógica de negocio off-chain en un único proceso. No utiliza NestJS; la arquitectura es un servidor Express monolítico (\texttt{src/server.ts}) con un proceso de indexación en segundo plano (\texttt{src/indexer.ts}).

\subsubsection{Servidor Express (\texttt{src/server.ts})}

Implementa todos los endpoints REST en un único archivo. Se organiza conceptualmente en los siguientes grupos de rutas:

\begin{itemize}
  \item \textbf{Auth}: \texttt{POST /api/auth/login}, \texttt{POST /api/auth/register}. Contraseñas con hash bcrypt (10 rounds), tokens JWT (7 días de expiración). Si \texttt{NODE\_ENV=production} y no se configura \texttt{JWT\_SECRET}, el servidor rechaza arrancar.
  \item \textbf{Usuarios}: \texttt{GET /api/users/me}, \texttt{PATCH /api/users/me}, \texttt{PATCH /api/users/me/wallet}. Roles: \texttt{CUSTOMER}, \texttt{ADMIN}, \texttt{STAFF}.
  \item \textbf{Eventos}: listado con caché en memoria (60 s), detalle por slug (30 s), destacados, tipos de boleto, eventos relacionados.
  \item \textbf{Carrito}: CRUD de ítems con verificación de propiedad por usuario (un usuario no puede modificar el carrito de otro).
  \item \textbf{Checkout}: previsualización y confirmación atómica (crea orden, ítems, boletos y pago en una sola transacción Prisma). Pago simulado como \texttt{PAID} (prototipo académico).
  \item \textbf{Órdenes y boletos}: historial de compras, boletos activos con campos Web3, boletos vendidos en P2P.
  \item \textbf{Transacciones Soroban}: \texttt{secure-ticket} (organizer-signed \texttt{crear\_boleto}), \texttt{list-ticket} (\texttt{listar\_boleto}), \texttt{cancel-listing} (\texttt{cancelar\_venta}), \texttt{build-buy-xdr} (construye XDR sin firmar para Freighter), \texttt{submit} (envía XDR firmado a la red).
  \item \textbf{Admin} (ADMIN): gestión de eventos, venues, contratos desplegados y escaneo QR. Protegido por verificación del rol en cada endpoint (re-consulta la BD).
  \item \textbf{Admin scan} (ADMIN/STAFF): \texttt{POST /api/admin/scan} recibe un \texttt{ticketId}, valida que el boleto exista y esté activo, y lo marca como \texttt{CANCELLED} en PostgreSQL. Nota de diseño: el scan es \textbf{DB-only} (sin llamada on-chain a \texttt{redimir\_boleto}) para minimizar la latencia en puerta; se trata de un trade-off documentado.
\end{itemize}

\subsubsection{Indexador Soroban (\texttt{src/indexer.ts})}

Proceso de fondo que sondea el RPC de Soroban cada 5 segundos. Procesa 7 tipos de eventos on-chain y actualiza PostgreSQL para mantener consistencia. Persiste el cursor en la tabla \texttt{indexer\_state} y se recupera automáticamente si el cursor cae fuera del rango de ledgers retenidos por el RPC público. No se ejecuta en entornos serverless (\texttt{VERCEL=1}).

\subsection{Capa de datos (Off-chain)}

La capa de datos off-chain se compone de una base de datos PostgreSQL 16 alojada en \textbf{Supabase}, gestionada mediante Prisma ORM. Utiliza un único schema \texttt{ticketing}:

\begin{itemize}
  \item \textbf{\texttt{users}}: id, email (único), password\_hash (bcrypt), first\_name, last\_name, phone, document\_type (CC/CE/TI/PP), document\_number, role (\texttt{CUSTOMER} $|$ \texttt{ADMIN} $|$ \texttt{STAFF}), \textbf{\texttt{wallet\_address}} (dirección Freighter, única, nullable).
  \item \textbf{\texttt{events}}: id, title, slug, status (DRAFT/PUBLISHED/CANCELLED/SOLD\_OUT), starts\_at, ends\_at, organizer\_id, venue\_id, category\_id, \textbf{\texttt{contract\_address}} (dirección del \texttt{ContratoEvento} desplegado en Testnet, nullable).
  \item \textbf{\texttt{tickets}}: id, order\_item\_id (FK), status (ACTIVE/USED/CANCELLED/REFUNDED), \textbf{\texttt{contract\_address}}, \textbf{\texttt{ticket\_root\_id}} (u32, ID estable on-chain), \textbf{\texttt{version}} (u32, se incrementa en cada reventa), \textbf{\texttt{is\_for\_sale}} (bool), \textbf{\texttt{resale\_price}} (BigInt, stroops), \textbf{\texttt{owner\_wallet}} (dirección del propietario en Stellar). Restricción única: \texttt{(contract\_address, ticket\_root\_id, version)}.
  \item \textbf{\texttt{indexer\_state}}: cursor de sincronización del indexador (\texttt{last\_ledger}).
  \item \textbf{Comercio Web2}: \texttt{organizers}, \texttt{venues}, \texttt{venue\_sections}, \texttt{seats}, \texttt{event\_ticket\_types}, \texttt{event\_seat\_inventory} (AVAILABLE/HELD/SOLD/BLOCKED), \texttt{carts}, \texttt{cart\_items}, \texttt{orders}, \texttt{order\_items}, \texttt{payments} (\texttt{provider\_reference} almacena el hash de transacción Soroban), \texttt{event\_categories}, \texttt{cities}.
\end{itemize}

PostgreSQL es la fuente de datos para consultas del frontend. Los datos on-chain se sincronizan en tiempo real mediante el Indexador Soroban que actualiza los campos Web3 de la tabla \texttt{tickets}.

\subsection{Capa blockchain (On-chain)}

La capa on-chain es la fuente de verdad del sistema para la propiedad de boletos y la ejecución de pagos. Todo el código de los contratos está escrito en español (nombres de structs, funciones, variables y errores). Se compone de:

\subsubsection{Contrato Factory (\texttt{FabricaBoletos})}

Contrato Soroban responsable de desplegar un contrato independiente por cada evento. Implementa el patrón Factory para garantizar aislamiento de configuración, storage y fondos entre eventos.

\paragraph{Modelo de datos.} Claves de storage (\texttt{ClaveDato}): \texttt{Administrador}, \texttt{ContadorEventos}, \texttt{HashWasmEvento} (hash del WASM compilado del \texttt{ContratoEvento}), \texttt{ContratoEvento(u32)} (mapeo \texttt{id\_evento} $\rightarrow$ dirección), \texttt{ContratoRegistrado(Address)} (mapeo inverso para evitar duplicados).

\paragraph{Configuración de evento.} El struct \texttt{ConfiguracionEvento} empaqueta: \texttt{id\_evento}, \texttt{organizador}, \texttt{token\_pago}, \texttt{comision\_organizador} (u32, 0--99), \texttt{comision\_plataforma} (u32, 0--99), \texttt{wallet\_organizador}, \texttt{wallet\_plataforma}, \texttt{capacidad\_total}.

\paragraph{Funciones públicas:}
\begin{itemize}
  \item \texttt{inicializar(administrador)}: setup único de la factory. Requiere auth del administrador.
  \item \texttt{configurar\_wasm\_evento(hash)}: almacena el hash WASM del \texttt{ContratoEvento}. Requiere auth del administrador.
  \item \texttt{crear\_evento\_contrato(config, dir\_prueba)} $\rightarrow$ \texttt{Address}: despliega e inicializa un nuevo contrato de evento. Requiere \textbf{doble autorización}: administrador y organizador. Valida comisiones $<$ 100\%, capacidad $>$ 0, id\_evento no duplicado.
  \item \texttt{obtener\_contrato\_evento(id\_evento)} $\rightarrow$ \texttt{Address}: consulta la dirección del contrato de un evento.
  \item \texttt{obtener\_contador\_eventos()} $\rightarrow$ \texttt{u32}: total de eventos creados.
\end{itemize}

\paragraph{Deploy programático.} En producción, usa \texttt{deployer().with\_current\_contract(salt).deploy\_v2(hash, ())} para crear contratos en la blockchain con dirección determinista derivada del \texttt{id\_evento}. En tests, usa una dirección pre-registrada (Soroban testutils no soporta deploy real).

\subsubsection{Contrato por Evento (\texttt{ContratoEvento})}

Contrato Soroban que implementa toda la lógica de negocio on-chain para un evento específico. Su estructura interna se organiza de la siguiente forma:

\paragraph{Modelo de datos (Storage).} El contrato utiliza un almacenamiento basado en claves (\texttt{ClaveDato}) con los siguientes elementos:

\begin{itemize}
  \item \texttt{Boleto(u32, u32)}: almacena el struct \texttt{Boleto} por \texttt{(ticket\_root\_id, version)}. Cada boleto contiene: \texttt{ticket\_root\_id} (ID estable que no cambia entre reventas), \texttt{version} (se incrementa con cada reventa: 0, 1, 2\ldots), \texttt{id\_evento}, \texttt{propietario} (Address), \texttt{precio} (i128), \texttt{en\_venta}, \texttt{es\_reventa}, \texttt{usado}, \texttt{invalidado}.
  \item \texttt{VersionActual(u32)}: puntero a la versión vigente de un boleto. Se actualiza en cada reventa.
  \item \texttt{ContadorBoletos}: contador autoincremental que asigna el próximo \texttt{ticket\_root\_id}.
  \item \texttt{Organizador}, \texttt{Plataforma}: direcciones de la ticketera y de la plataforma.
  \item \texttt{TokenPago}: dirección del SAC XLM nativo (\texttt{CDLZFC3SYJYDZT7K67VZ75HPJVIEUVNIXF47ZG2FB2RMQQVU2HHGCYSC}).
  \item \texttt{ComisionOrganizador}, \texttt{ComisionPlataforma}: porcentajes de comisión (i128, base 100).
  \item \texttt{Verificador(Address)}: registra si una dirección tiene permiso para redimir boletos en puerta.
\end{itemize}

\paragraph{Funciones de escritura (state-changing).} Cada función está protegida por \texttt{require\_auth()} que verifica la firma del actor autorizado:

\begin{itemize}
  \item \texttt{inicializar}: configura el contrato con las direcciones, token de pago y porcentajes de comisión. Solo puede ejecutarse una vez (guarda de re-inicialización). Requiere auth del organizador.
  \item \texttt{crear\_boleto(id\_evento, precio)} $\rightarrow$ \texttt{u32}: emite un nuevo boleto (version 0) asignado al organizador. Requiere auth del organizador. Emite evento \texttt{BoletoCreado}.
  \item \texttt{listar\_boleto(ticket\_root\_id, nuevo\_precio)}: marca un boleto como en venta y actualiza su precio. Requiere auth del propietario. Valida que no esté ya en venta, ni usado, ni invalidado. Emite evento \texttt{BoletoListado}.
  \item \texttt{cancelar\_venta(ticket\_root\_id)}: retira un boleto del marketplace. Requiere auth del propietario. Emite evento \texttt{VentaCancelada}.
  \item \texttt{comprar\_boleto(ticket\_root\_id, comprador)} $\rightarrow$ \texttt{u32}: ejecuta la compra. Implementa dos flujos en una sola transacción atómica:
    \begin{itemize}
      \item \textit{Venta primaria} (\texttt{es\_reventa = false}): el 100\% del pago se transfiere al organizador mediante \texttt{ClienteToken}. El boleto se actualiza in-place: nuevo propietario, \texttt{es\_reventa = true}. Retorna la versión actual. Emite evento \texttt{BoletoCompradoPrimario}.
      \item \textit{Reventa} (\texttt{es\_reventa = true}): el pago se distribuye en tres transferencias atómicas: comisión al organizador, comisión a la plataforma y remanente al vendedor. Luego se ejecuta \textbf{burn/remint}: la versión actual se marca como \texttt{invalidado = true}, se crea una nueva versión (\texttt{version + 1}) asignada al comprador, y se actualiza el puntero \texttt{VersionActual}. Retorna la nueva versión. Emite evento \texttt{BoletoRevendido}.
    \end{itemize}
    Requiere auth del comprador. Valida que el comprador no sea el propietario actual (\texttt{AutoCompra}).
  \item \texttt{agregar\_verificador(address)}: registra un verificador autorizado. Requiere auth del organizador. Emite evento \texttt{VerificadorAgregado}.
  \item \texttt{remover\_verificador(address)}: desautoriza un verificador. Requiere auth del organizador. Emite evento \texttt{VerificadorRemovido}.
  \item \texttt{redimir\_boleto(ticket\_root\_id, verificador)}: marca el boleto como \texttt{usado = true} (irreversible). \textbf{Solo verificadores autorizados} pueden invocarla (no el propietario). Requiere auth del verificador. Emite evento \texttt{BoletoRedimido}.
  \item \texttt{invalidar\_boleto(ticket\_root\_id)}: cancela un boleto administrativamente (\texttt{invalidado = true}). Requiere auth del organizador. Emite evento \texttt{BoletoInvalidadoEvt}.
\end{itemize}

\paragraph{Funciones de lectura (views).} No requieren autorización:

\begin{itemize}
  \item \texttt{obtener\_boleto(ticket\_root\_id)}: retorna el struct \texttt{Boleto} de la versión vigente.
  \item \texttt{obtener\_boleto\_version(ticket\_root\_id, version)}: retorna una versión específica (historial).
  \item \texttt{obtener\_version\_vigente(ticket\_root\_id)}: retorna el número de versión actual.
  \item \texttt{obtener\_propietario(ticket\_root\_id)}: retorna la dirección del propietario actual.
  \item \texttt{obtener\_boletos\_reventa()}: retorna todos los boletos en reventa activa (O(n), usar indexador off-chain en producción).
  \item \texttt{obtener\_boletos\_evento(id\_evento)}: retorna todos los boletos de un evento (O(n)).
  \item \texttt{es\_verificador(address)}: consulta si una dirección es verificador autorizado.
\end{itemize}

\paragraph{Errores tipados (\texttt{ErrorContrato}).} El contrato define 16 errores tipados como enum con \texttt{\#[contracterror]} y \texttt{\#[repr(u32)]}, lo que permite manejo programático off-chain sin parsear strings:

\begin{table}[htbp]
  \centering
  \small
  \begin{tabularx}{\textwidth}{@{}c l X@{}}
    \toprule
    \textbf{Código} & \textbf{Nombre} & \textbf{Significado}\\
    \midrule
    1 & YaInicializado & Doble inicialización del contrato\\
    2 & ComisionesMuyAltas & Suma de comisiones $\geq$ 100\%\\
    3 & ComisionesNegativas & Comisión menor a 0\\
    4 & PrecioInvalido & Precio $\leq$ 0\\
    5 & YaEnVenta & Boleto ya listado\\
    6 & BoletoUsado & Boleto ya redimido\\
    7 & NoEnVenta & Boleto no está en venta\\
    8 & AutoCompra & Comprador = propietario\\
    9 & YaUsado & Segundo intento de redención\\
    10 & BoletoNoEncontrado & ID o versión inexistente\\
    11 & NoAutorizado & Sin permisos para la acción\\
    12 & VersionInvalida & Versión no existe\\
    13 & BoletoInvalidado & Boleto cancelado/burned\\
    14 & NoInicializado & Contrato sin inicializar\\
    15 & VerificadorYaExiste & Dirección ya es verificador\\
    16 & VerificadorNoEncontrado & Dirección no es verificador\\
    \bottomrule
  \end{tabularx}
  \caption{Errores tipados del \texttt{ContratoEvento}}
  \label{tab:errores-contrato}
\end{table}

\paragraph{Eventos on-chain.} El contrato emite eventos tipados con \texttt{\#[contractevent]} para indexación off-chain. Cada evento tiene hasta 4 topics indexados:

\texttt{BoletoCreado}, \texttt{BoletoListado}, \texttt{VentaCancelada}, \texttt{BoletoCompradoPrimario}, \texttt{BoletoRevendido}, \texttt{BoletoRedimido}, \texttt{BoletoInvalidadoEvt}, \texttt{VerificadorAgregado}, \texttt{VerificadorRemovido}.

\subsubsection{SAC XLM Nativo (Stellar Asset Contract)}

Contrato externo que implementa el estándar de token de Stellar para el XLM nativo. El contrato de evento lo invoca mediante \texttt{ClienteToken} (alias de \texttt{soroban\_sdk::token::Client}) para ejecutar las transferencias de pago. Se usa el SAC del XLM nativo, cuya dirección en testnet es \texttt{CDLZFC3S...HHGCYSC}. Los precios se expresan en \textit{stroops} (1 XLM = 10\textsuperscript{7} stroops).

\subsection{Separación on-chain / off-chain}

La siguiente tabla resume el criterio de distribución de responsabilidades entre ambas capas:

\begin{table}[htbp]
  \centering
  \begin{tabularx}{\textwidth}{@{}X l l@{}}
    \toprule
    \textbf{Responsabilidad} & \textbf{On-chain} & \textbf{Off-chain}\\
    \midrule
    Propiedad de boletos (quién es dueño) & \checkmark & Copia indexada\\
    Ejecución de pagos y comisiones & \checkmark & ---\\
    Creación de boletos (mint) & \checkmark & Registro en PostgreSQL\\
    Redención (marcar como usado) & \checkmark & Copia indexada\\
    Invalidación administrativa & \checkmark & Copia indexada\\
    Gestión de verificadores autorizados & \checkmark & ---\\
    Historial de versiones (burn/remint) & \checkmark & Consulta indexada\\
    Despliegue de contratos por evento & \checkmark (Factory) & ---\\
    Consultas de búsqueda y filtrado & --- & \checkmark\\
    Autenticación y roles de usuario & --- & \checkmark\\
    Carrito y seat holds & --- & \checkmark\\
    Gestión de eventos, sedes y asientos & --- & \checkmark\\
    \bottomrule
  \end{tabularx}
  \caption{Distribución de responsabilidades on-chain vs. off-chain}
  \label{tab:onchain-offchain}
\end{table}

El principio rector es: \textit{lo que afecta propiedad o dinero se ejecuta on-chain} (fuente de verdad inmutable), \textit{lo que requiere consultas rápidas o contexto de negocio se resuelve off-chain} (rendimiento y usabilidad).

%===============================================================
\FloatBarrier
\clearpage
\section{Vista de procesos}
\label{sec:vista-procesos}
%===============================================================

Esta sección describe los flujos de interacción principales del sistema, enfocándose en cómo los componentes colaboran en tiempo de ejecución para resolver los casos de uso críticos. Se documentan tres flujos y la máquina de estados que gobierna el ciclo de vida del ticket.

% ---- Flujo 1: Compra primaria ----
\subsection{Flujo de compra primaria}

La compra primaria es el primer flujo que ocurre sobre un boleto recién creado. Involucra al comprador, el frontend TuTicket, el Backend API, el contrato de evento y el token de pago.

\begin{enumerate}
  \item El organizador ha creado previamente el evento y emitido boletos mediante \texttt{crear\_boleto}. Cada boleto queda asignado al organizador con \texttt{en\_venta = false} y \texttt{es\_reventa = false}.
  \item El organizador lista el boleto para venta primaria invocando \texttt{listar\_boleto}, que establece \texttt{en\_venta = true} y el precio.
  \item El comprador visualiza el boleto en la aplicación web TuTicket y selecciona la compra.
  \item El frontend envía la solicitud al Backend API, que orquesta la invocación de \texttt{comprar\_boleto} en el contrato de evento.
  \item El contrato ejecuta la lógica de compra primaria: como \texttt{es\_reventa = false}, el 100\% del pago se transfiere al organizador mediante \texttt{ClienteToken}.
  \item El contrato actualiza el boleto: \texttt{propietario = comprador}, \texttt{en\_venta = false}, \texttt{es\_reventa = true} (todas las ventas futuras serán reventas). Emite evento \texttt{BoletoCompradoPrimario}.
  \item El backend registra la operación en PostgreSQL.
\end{enumerate}

\textbf{Aspecto arquitectónico clave:} la compra primaria establece la transición de \texttt{es\_reventa = false} a \texttt{true}, lo que determina el comportamiento de todas las transacciones futuras sobre ese boleto (activación del split de comisiones y del mecanismo burn/remint).

% ---- Flujo 2: Reventa atómica ----
\subsection{Flujo de reventa atómica}

La Figura~\ref{fig:seq-reventa} presenta el diagrama de secuencia completo de este flujo.

\begin{figure}[htbp]
  \centering
  \includegraphics[width=0.8\textwidth]{04_secuencia_reventa_atomica.png}
  \caption{Diagrama de secuencia -- Reventa atómica con distribución de comisiones}
  \label{fig:seq-reventa}
\end{figure}

Este es el flujo más crítico del sistema desde el punto de vista arquitectónico. La reventa debe ser \textbf{atómica}: todas las operaciones (distribución de pagos, burn de la versión anterior y remint de nueva versión) ocurren en una sola transacción de Soroban, o ninguna ocurre. No existe estado intermedio.

\begin{enumerate}
  \item \textbf{Publicación:} el vendedor (propietario actual) invoca \texttt{listar\_boleto} con el nuevo precio. El contrato verifica que el vendedor es el propietario (\texttt{require\_auth}), que el boleto no está ya en venta, ni usado, ni invalidado. El boleto queda con \texttt{en\_venta = true}. El contrato emite evento \texttt{BoletoListado}.
  \item \textbf{Compra:} el comprador selecciona el boleto en el marketplace e inicia la compra. El frontend envía la solicitud al Backend API, que orquesta la invocación de \texttt{comprar\_boleto}.
  \item \textbf{Ejecución atómica en el contrato:} como \texttt{es\_reventa = true}, el contrato ejecuta la distribución de comisiones y el burn/remint en una sola transacción:
    \begin{enumerate}[label=\alph*)]
      \item Calcula la comisión del organizador: \texttt{precio $\times$ ComisionOrganizador / 100}.
      \item Calcula la comisión de la plataforma: \texttt{precio $\times$ ComisionPlataforma / 100}.
      \item Calcula el pago neto al vendedor: \texttt{precio $-$ comisión\_org $-$ comisión\_plat}.
      \item Ejecuta tres transferencias mediante \texttt{ClienteToken}: al organizador, a la plataforma y al vendedor.
      \item \textbf{Burn:} marca la versión actual del boleto como \texttt{invalidado = true}, \texttt{en\_venta = false}.
      \item \textbf{Remint:} crea nueva versión (\texttt{version + 1}) con \texttt{propietario = comprador}, \texttt{en\_venta = false}, \texttt{es\_reventa = true}.
      \item Actualiza el puntero \texttt{VersionActual} a la nueva versión.
    \end{enumerate}
  \item \textbf{Evento on-chain:} el contrato emite \texttt{BoletoRevendido} con \texttt{ticket\_root\_id}, \texttt{vendedor}, \texttt{comprador}, \texttt{precio}, \texttt{version\_anterior} y \texttt{version\_nueva}.
  \item \textbf{Registro off-chain:} el backend registra la transacción en PostgreSQL.
  \item \textbf{Notificación:} el frontend muestra al comprador su nuevo boleto y notifica al vendedor que la venta se completó.
\end{enumerate}

\textbf{Garantía de atomicidad:} si cualquiera de las transferencias de pago falla (por ejemplo, fondos insuficientes del comprador), toda la transacción se revierte automáticamente por el runtime de Soroban. Esto elimina por diseño los estados inconsistentes donde el pago se procesa pero el boleto no se transfiere, o viceversa. No se requiere lógica de compensación ni rollback manual.

\textbf{Trazabilidad por versionado:} la combinación \texttt{(ticket\_root\_id, version)} permite reconstruir el historial completo de propiedad de un boleto consultando todas sus versiones. Las versiones anteriores quedan como registros históricos inmutables con \texttt{invalidado = true}.

% ---- Flujo 3: Verificación y redención en puerta ----
\subsection{Flujo de verificación y redención en puerta}

La verificación de acceso al evento involucra a un verificador autorizado que valida el boleto del asistente y lo marca como redimido on-chain.

\begin{enumerate}
  \item El organizador ha creado previamente una cuenta con rol \texttt{STAFF} (o \texttt{ADMIN}) en el sistema, asignada al personal de verificación.
  \item El asistente presenta su boleto (código QR) en la puerta del evento. El QR codifica un payload JSON con el \texttt{ticketId} (UUID del boleto en PostgreSQL).
  \item El verificador utiliza la página \texttt{/escanear} del sistema (accesible desde el navegador del dispositivo de verificación) con su sesión activa.
  \item La cámara del dispositivo lee el QR. El frontend envía \texttt{POST /api/admin/scan \{ticketId\}} al Backend API.
  \item El backend valida: el boleto existe en PostgreSQL, tiene estado \texttt{ACTIVE}. Si no, retorna error con mensaje descriptivo.
  \item Si es válido, el backend actualiza el estado a \texttt{CANCELLED} en PostgreSQL y retorna \texttt{\{success: true\}}.
  \item La interfaz del escáner muestra ``Acceso Permitido'' (verde) o ``Acceso Denegado'' (rojo) con el código del boleto.
\end{enumerate}

\textbf{Nota de implementación:} el escáner actual opera en modo \textbf{DB-only} (sin llamada on-chain a \texttt{redimir\_boleto}). Esta decisión de diseño minimiza la latencia en puerta (evita el round-trip al RPC Soroban de \textasciitilde5 segundos). El contrato Soroban implementa la función \texttt{redimir\_boleto} y el sistema de verificadores autorizados on-chain como mecanismo de validación criptográfica, que puede activarse en una versión de producción donde la latencia adicional sea aceptable.

\textbf{Aspecto arquitectónico clave:} el sistema de verificadores on-chain (registro en el contrato Soroban) y la gestión de roles off-chain (STAFF en PostgreSQL) son mecanismos complementarios que refuerzan la separación entre propietario (custodia del ticket) y personal autorizado (capacidad de redención).

La Figura~\ref{fig:seq-verificacion} presenta el diagrama de secuencia completo de este flujo.

\begin{figure}[htbp]
  \centering
  \includegraphics[width=0.85\textwidth]{05_secuencia_verificacion_puerta.png}
  \caption{Diagrama de secuencia -- Verificación en puerta con escáner QR (Phase 4)}
  \label{fig:seq-verificacion}
\end{figure}

% ---- Máquina de estados ----
\subsection{Máquina de estados del ticket}

El ciclo de vida de un ticket está gobernado por una máquina de estados que define las transiciones válidas entre estados y las condiciones bajo las cuales pueden ocurrir. Esta máquina se implementa implícitamente en el contrato Soroban mediante las validaciones en cada función. La Figura~\ref{fig:estados-ticket} presenta el diagrama de estados completo.

\begin{figure}[htbp]
  \centering
  \includegraphics[width=0.5\textwidth]{06_estados_ticket.png}
  \caption{Máquina de estados del ciclo de vida de un ticket}
  \label{fig:estados-ticket}
\end{figure}

\subsubsection{Estados}

\begin{table}[h!]
  \centering
  \begin{tabularx}{\textwidth}{@{}p{3.5cm} X@{}}
    \toprule
    \textbf{Estado} & \textbf{Descripción}\\
    \midrule
    \textbf{Emitido} & El boleto acaba de ser creado por el organizador (\texttt{crear\_boleto}). Está asignado al organizador, con \texttt{en\_venta = false}, \texttt{es\_reventa = false}, \texttt{usado = false}, \texttt{invalidado = false}.\\
    \textbf{En venta} & El boleto ha sido listado para venta o reventa (\texttt{listar\_boleto}). \texttt{en\_venta = true}. Visible en el marketplace.\\
    \textbf{Asignado} & El boleto tiene un propietario (comprador) y no está en venta. \texttt{en\_venta = false}. El propietario puede listarlo para reventa o presentarlo en la puerta del evento.\\
    \textbf{Consumido} & El boleto fue redimido por un verificador autorizado (\texttt{redimir\_boleto}). \texttt{usado = true}. Estado terminal e irreversible. No puede listarse, comprarse ni revenderse.\\
    \textbf{Invalidado} & El boleto fue cancelado administrativamente por el organizador (\texttt{invalidar\_boleto}) o fue ``quemado'' (\textit{burned}) durante una reventa para crear la nueva versión. \texttt{invalidado = true}. Estado terminal on-chain.\\
    \bottomrule
  \end{tabularx}
  \caption{Estados del ciclo de vida de un boleto}
  \label{tab:estados-ticket}
\end{table}

\subsubsection{Transiciones}

\begin{table}[htbp]
  \centering
  \begin{tabularx}{\textwidth}{@{}p{2.5cm} p{2.5cm} p{3.5cm} X@{}}
    \toprule
    \textbf{Desde} & \textbf{Hacia} & \textbf{Acción} & \textbf{Condiciones}\\
    \midrule
    Emitido & En venta & \texttt{listar\_boleto} & El organizador (propietario) lista el boleto para venta primaria.\\
    En venta & Asignado & \texttt{comprar\_boleto} & Venta primaria: 100\% al organizador, \texttt{es\_reventa} pasa a \texttt{true}. Update in-place.\\
    Asignado & En venta & \texttt{listar\_boleto} & El propietario lista el boleto para reventa. No debe estar usado ni invalidado.\\
    En venta & Asignado & \texttt{cancelar\_venta} & El propietario retira el boleto del marketplace.\\
    En venta & Asignado (nueva versión) & \texttt{comprar\_boleto} & Reventa: distribución de comisiones + burn versión actual + remint nueva versión.\\
    Asignado & Consumido & \texttt{redimir\_boleto} & Un verificador autorizado redime el boleto. Irreversible (\texttt{usado = true}).\\
    No invalidado & Invalidado & \texttt{invalidar\_boleto} & El organizador cancela el boleto on-chain. Irreversible (\texttt{invalidado = true}).\\
    En venta (reventa) & Invalidado & \texttt{comprar\_boleto} & La versión previa se invalida (burn) como parte del flujo de reventa.\\
    \bottomrule
  \end{tabularx}
  \caption{Transiciones válidas entre estados del boleto}
  \label{tab:transiciones-ticket}
\end{table}

\subsubsection{Reglas de integridad}

Las siguientes reglas son impuestas por el contrato y no pueden ser violadas:

\begin{itemize}
  \item Un boleto con \texttt{usado = true} no puede cambiar a ningún otro estado. La redención es \textbf{irreversible}.
  \item Un boleto con \texttt{invalidado = true} no puede venderse, comprarse, listarse ni redimirse. La invalidación es \textbf{irreversible}.
  \item Solo el \texttt{propietario} actual puede invocar \texttt{listar\_boleto} y \texttt{cancelar\_venta} (verificado por \texttt{require\_auth}).
  \item Solo un \textbf{verificador autorizado} (registrado on-chain) puede invocar \texttt{redimir\_boleto}. El propietario no puede auto-redimirse.
  \item Un propietario no puede comprarse su propio boleto (\texttt{AutoCompra} está prohibido).
  \item El contrato solo puede inicializarse una vez (error \texttt{YaInicializado}).
  \item Las comisiones de organizador y plataforma no pueden sumar 100\% o más (\texttt{ComisionesMuyAltas}).
  \item Cada versión de un boleto es inmutable una vez invalidada: las versiones anteriores quedan como registro histórico en el storage.
\end{itemize}

%===============================================================
\FloatBarrier
\clearpage
\section{Vista de despliegue}
\label{sec:vista-despliegue}
%===============================================================

El sistema se despliega en cuatro servicios cloud independientes que conforman la topología de producción.

\begin{table}[htbp]
\centering
\caption{Nodos de despliegue}
\begin{tabularx}{\textwidth}{l l X}
\toprule
\textbf{Nodo} & \textbf{Plataforma} & \textbf{Componentes} \\
\midrule
Cliente & Navegador (Chrome + Freighter) & SPA React servida como archivos estáticos. Extensión Freighter gestiona clave privada del usuario. \\
Frontend & Vercel (CDN estático) & Bundle \texttt{/frontend/dist}. SPA con 23 rutas client-side. \\
Backend + Indexador & Railway / Render (proceso long-lived) & \texttt{tsx src/server.ts}: API Express + Indexador Soroban en background. El Indexador requiere proceso persistente (no compatible con serverless). \\
Base de datos & Supabase (PostgreSQL 16) & Schema \texttt{ticketing}: 26+ tablas Web2+Web3. Pool de conexiones vía pgbouncer (\texttt{connection\_limit=5}). \\
Blockchain & Stellar Testnet (Soroban) & \texttt{FabricaBoletos} + 11 instancias de \texttt{ContratoEvento} + SAC XLM nativo. \\
\bottomrule
\end{tabularx}
\end{table}

\subsection{Configuración de entornos}

El sistema distingue dos modos de ejecución para el backend:

\begin{itemize}
  \item \textbf{Long-lived (Railway/Render)}: ejecuta \texttt{tsx src/server.ts}. Arranca \texttt{app.listen(PORT)} e inicia el Indexador Soroban en segundo plano.
  \item \textbf{Serverless (Vercel)}: la variable \texttt{VERCEL=1} desactiva \texttt{app.listen()} y el Indexador. El backend se exporta como handler via \texttt{backend/api/index.ts}.
\end{itemize}

En desarrollo local, el frontend corre con \texttt{npm run dev} (Vite, puerto 8080) y el backend con \texttt{npm run dev} (tsx watch, puerto 3000).

\subsection{Variables de entorno relevantes}

\begin{table}[htbp]
\centering
\caption{Variables de entorno principales}
\begin{tabularx}{\textwidth}{l l X}
\toprule
\textbf{Variable} & \textbf{Servicio} & \textbf{Descripción} \\
\midrule
\texttt{DATABASE\_URL} & Backend & Cadena de conexión PostgreSQL con Supabase (Prisma, \texttt{connection\_limit=5}). \\
\texttt{JWT\_SECRET} & Backend & Secreto para firma de tokens JWT. Obligatorio en producción; falla al arrancar si no está definido. \\
\texttt{SOROBAN\_RPC\_URL} & Backend & URL del nodo RPC de Soroban (\texttt{https://soroban-testnet.stellar.org}). \\
\texttt{ORGANIZER\_SECRET} & Backend & Clave privada del organizador para firmar server-side \texttt{crear\_boleto}, \texttt{listar\_boleto} y \texttt{cancelar\_venta}. \\
\texttt{FACTORY\_CONTRACT\_ID} & Backend & ID del contrato \texttt{FabricaBoletos} desplegado. Requerido por el panel de administración. \\
\texttt{VITE\_API\_BASE\_URL} & Frontend & URL base del backend API (por defecto \texttt{http://localhost:3000}). \\
\bottomrule
\end{tabularx}
\end{table}

\begin{figure}[htbp]
  \centering
  \includegraphics[width=0.95\textwidth]{03_despliegue_produccion.png}
  \caption{Diagrama de despliegue -- Topología de nodos en producción (Vercel + Railway/Render + Supabase + Stellar Testnet)}
  \label{fig:despliegue}
\end{figure}

%===============================================================
\FloatBarrier
\clearpage
\section{Vista de implementación}
\label{sec:vista-implementacion}
%===============================================================

\subsection{Organización de repositorios}

El proyecto se organiza en un \textbf{monorepo} principal que contiene contratos, backend y frontend, más un repositorio separado para la documentación:

\begin{table}[htbp]
\centering
\caption{Repositorios del proyecto}
\begin{tabularx}{\textwidth}{l l X}
\toprule
\textbf{Repositorio} & \textbf{Ruta local} & \textbf{Contenido} \\
\midrule
\texttt{main\_contract} & \texttt{codigo/main\_contract/} & \textbf{Monorepo}: contratos Soroban (Rust), backend Express (Node.js) y frontend React. \\
\texttt{Documento\_Tesis} & \texttt{Documentos/Documento\_Tesis/} & Documentación LaTeX del proyecto. \\
\bottomrule
\end{tabularx}
\end{table}

La estructura interna del monorepo \texttt{main\_contract} es:

\begin{verbatim}
main_contract/
+-- contracts/     # Cargo workspace: event_contract + factory_contract
+-- backend/       # Node.js Express + Prisma + Indexador Soroban
+-- frontend/      # React 18 + Vite + TypeScript + Freighter
+-- docs/          # Documentación técnica (architecture, ADRs, backlog)
+-- .github/       # CI: build contratos + tests cargo
\end{verbatim}

\subsection{Estructura del workspace de contratos}

El repositorio de contratos utiliza un workspace de Cargo con dos crates:

\begin{verbatim}
contracts/
+-- Cargo.toml              # Workspace raiz
+-- contracts/
    +-- event_contract/
    |   +-- src/
    |       +-- lib.rs       # ContratoEvento (~1140 lineas)
    |       +-- test.rs      # Suite de tests
    +-- factory_contract/
        +-- src/
            +-- lib.rs       # FabricaBoletos (~430 lineas)
            +-- test.rs      # Suite de tests
\end{verbatim}

Dependencia: \texttt{factory\_contract} depende de \texttt{event\_contract} (importa \texttt{ContratoEventoClient} para inicializar contratos hijos).

\subsection{Estructura del backend}

\begin{verbatim}
src/
+-- app.module.ts            # Modulo raiz
+-- auth/                    # Autenticacion JWT (registro, login, guards)
+-- users/                   # Gestion de usuarios y perfiles
+-- events/                  # CRUD de eventos
+-- tickets/                 # Gestion de boletos y disponibilidad
+-- cart/                    # Carrito de compras (sesion)
+-- checkout/                # Procesamiento de compras (simulado)
+-- admin/                   # Panel de administracion
+-- prisma/                  # Servicio Prisma y migraciones
+-- common/                  # Guards, decoradores, DTOs compartidos
\end{verbatim}

\subsection{Estructura del frontend}

\begin{verbatim}
src/
+-- main.tsx                 # Entry point
+-- App.tsx                  # Router principal
+-- components/              # Componentes reutilizables (shadcn/ui)
+-- pages/                   # Vistas por ruta
+-- services/                # Llamadas HTTP al backend (Axios)
+-- context/                 # React Context (auth, carrito)
+-- hooks/                   # Custom hooks
+-- types/                   # Interfaces TypeScript
\end{verbatim}

\subsection{Artefactos compilados y comandos de build}

\begin{table}[htbp]
\centering
\caption{Comandos de compilación y artefactos}
\small
\begin{tabularx}{\textwidth}{l l X}
\toprule
\textbf{Componente} & \textbf{Comando} & \textbf{Artefacto} \\
\midrule
Contratos & \texttt{stellar contract build} & \texttt{target/wasm32v1-none/release/\{event,factory\}\_contract.wasm} \\
Backend & \texttt{npm run build} & \texttt{dist/} (JavaScript transpilado) \\
Frontend & \texttt{npm run build} & \texttt{dist/} (SPA estática optimizada) \\
Tests contratos & \texttt{cargo test} & Resultados de pruebas unitarias Soroban (48+ tests) \\
\bottomrule
\end{tabularx}
\end{table}

%===============================================================
\FloatBarrier
\clearpage
\section{Vista de datos}
\label{sec:vista-datos}
%===============================================================

El sistema mantiene dos modelos de datos complementarios: on-chain (inmutable, en los contratos Soroban) y off-chain (relacional, en PostgreSQL).

\subsection{Modelo de datos on-chain: ContratoEvento}

\subsubsection{Claves de storage (\texttt{ClaveDato})}

\begin{table}[htbp]
\centering
\caption{Claves de storage del \texttt{ContratoEvento}}
\begin{tabularx}{\textwidth}{l l X}
\toprule
\textbf{Clave} & \textbf{Tipo de valor} & \textbf{Descripción} \\
\midrule
\texttt{Organizador} & \texttt{Address} & Dirección del organizador del evento. \\
\texttt{Plataforma} & \texttt{Address} & Wallet de la plataforma para comisiones. \\
\texttt{TokenPago} & \texttt{Address} & Contrato del token usado para pagos. \\
\texttt{ComisionOrganizador} & \texttt{i128} & Porcentaje de comisión del organizador. \\
\texttt{ComisionPlataforma} & \texttt{i128} & Porcentaje de comisión de la plataforma. \\
\texttt{ContadorBoletos} & \texttt{u32} & ID incremental para nuevos boletos. \\
\texttt{Boleto(u32)} & \texttt{Boleto} & Datos de un boleto por su ID compuesto. \\
\texttt{VersionActual(u32)} & \texttt{u32} & Versión vigente de un \texttt{ticket\_root\_id}. \\
\texttt{EsVerificador(Address)} & \texttt{bool} & Indica si una dirección es verificador autorizado. \\
\bottomrule
\end{tabularx}
\end{table}

\subsubsection{Estructura del boleto (\texttt{Boleto})}

\begin{table}[htbp]
\centering
\caption{Campos del struct \texttt{Boleto}}
\begin{tabularx}{\textwidth}{l l X}
\toprule
\textbf{Campo} & \textbf{Tipo} & \textbf{Descripción} \\
\midrule
\texttt{ticket\_root\_id} & \texttt{u32} & Identificador raíz (invariante entre versiones). \\
\texttt{version} & \texttt{u32} & Número de versión (incrementa con cada reventa). \\
\texttt{id\_evento} & \texttt{u32} & ID del evento al que pertenece. \\
\texttt{propietario} & \texttt{Address} & Dirección del dueño actual. \\
\texttt{precio} & \texttt{i128} & Precio en unidades del token de pago. \\
\texttt{en\_venta} & \texttt{bool} & Si está listado para reventa. \\
\texttt{es\_reventa} & \texttt{bool} & \texttt{true} después de la primera compra. \\
\texttt{usado} & \texttt{bool} & Si fue redimido (irreversible). \\
\texttt{invalidado} & \texttt{bool} & Si fue invalidado por el organizador (irreversible). \\
\bottomrule
\end{tabularx}
\end{table}

\FloatBarrier
\subsection{Modelo de datos on-chain: FabricaBoletos}

El contrato \texttt{FabricaBoletos} mantiene un registro centralizado de todos los eventos desplegados y su configuración. Su storage es de sólo escritura para el administrador y de lectura pública para el backend.

\begin{table}[htbp]
\centering
\caption{Claves de storage de \texttt{FabricaBoletos}}
\begin{tabularx}{\textwidth}{l l X}
\toprule
\textbf{Clave} & \textbf{Tipo de valor} & \textbf{Descripción} \\
\midrule
\texttt{Administrador} & \texttt{Address} & Admin de la factory. \\
\texttt{ContadorEventos} & \texttt{u32} & Total de eventos creados. \\
\texttt{HashWasmEvento} & \texttt{BytesN<32>} & Hash del WASM del \texttt{ContratoEvento}. \\
\texttt{ContratoEvento(u32)} & \texttt{Address} & Mapeo id\_evento → dirección del contrato. \\
\texttt{ContratoRegistrado(Address)} & \texttt{bool} & Mapeo inverso para evitar duplicados. \\
\bottomrule
\end{tabularx}
\end{table}

\FloatBarrier
\subsection{Modelo de datos off-chain: PostgreSQL}

La base de datos relacional utiliza el schema \texttt{ticketing} en PostgreSQL 16 (hosteado en Supabase):

\subsubsection{Entidades principales}

El esquema contiene las siguientes tablas, con sus campos Web3 añadidos como capa de extensión sobre el modelo relacional clásico, con sus campos Web3 añadidos como capa de extensión sobre el modelo relacional clásico:

\begin{itemize}
  \item \textbf{\texttt{users}}: id, email (único), password\_hash (bcrypt), nombre, teléfono, tipo y número de documento, rol (\texttt{CUSTOMER} | \texttt{ADMIN} | \texttt{STAFF}), \textbf{\texttt{wallet\_address}} (dirección pública Freighter, único, opcional).
  \item \textbf{\texttt{events}}: id, título, slug (único), descripción, fecha, estado (\texttt{DRAFT} | \texttt{PUBLISHED} | \texttt{CANCELLED}), venue, organizador, categoría, \textbf{\texttt{contract\_address}} (dirección del \texttt{ContratoEvento} en Soroban, único, opcional).
  \item \textbf{\texttt{tickets}}: id, estado (\texttt{ACTIVE} | \texttt{USED} | \texttt{CANCELLED} | \texttt{REFUNDED}), order\_item\_id (FK), \textbf{\texttt{contract\_address}}, \textbf{\texttt{ticket\_root\_id}} (\texttt{u32}), \textbf{\texttt{version}} (\texttt{u32}), \textbf{\texttt{is\_for\_sale}} (bool), \textbf{\texttt{resale\_price}} (BIGINT en stroops), \textbf{\texttt{owner\_wallet}}. Restricción de unicidad: \texttt{(contract\_address, ticket\_root\_id, version)}.
  \item \textbf{\texttt{orders}} / \textbf{\texttt{order\_items}}: historial de compras fiat. subtotal, service\_fees, total.
  \item \textbf{\texttt{payments}}: método (\texttt{CARD} | \texttt{PSE} | \texttt{CASHPOINT}), estado, monto, \textbf{\texttt{provider\_reference}} (hash de tx Soroban cuando aplica).
  \item \textbf{\texttt{carts}} / \textbf{\texttt{cart\_items}}: carrito de compras por usuario.
  \item \textbf{\texttt{event\_ticket\_types}}: tipos de boleta por evento (nombre, precio, cantidad).
  \item \textbf{\texttt{indexer\_state}}: cursor del Indexador Soroban (\texttt{last\_ledger}), un único registro.
\end{itemize}

La Figura~\ref{fig:erd-datos} presenta el diagrama entidad-relación de las tablas clave con sus campos Web3.

\begin{figure}[htbp]
  \centering
  \includegraphics[width=\textwidth]{08_erd_datos.png}
  \caption{Diagrama ERD -- Modelo de datos híbrido Web2.5 con campos Web3}
  \label{fig:erd-datos}
\end{figure}

\subsubsection{Correspondencia on-chain / off-chain}

\begin{table}[htbp]
\centering
\caption{Correspondencia entre datos on-chain y off-chain}
\begin{tabularx}{\textwidth}{p{2.6cm} p{5.2cm} X}
\toprule
\textbf{Concepto} & \textbf{On-chain} & \textbf{Off-chain} \\
\midrule
Evento & \texttt{ContratoEvento} (contrato por evento) & Tabla \texttt{events} + campo \texttt{contract\_address} \\
Boleto & \texttt{Boleto} struct: \texttt{ticket\_root\_id}, \texttt{version}, \texttt{propietario} & Tabla \texttt{tickets} + campos \texttt{ticket\_root\_id}, \texttt{version}, \texttt{owner\_wallet} \\
Propietario & \texttt{propietario} (Address Stellar) & \texttt{owner\_wallet} en \texttt{tickets} + \texttt{wallet\_address} en \texttt{users} \\
En venta & \texttt{en\_venta} (bool) & \texttt{is\_for\_sale} (bool) + \texttt{resale\_price} (BIGINT, stroops) \\
Estado de uso & \texttt{usado} / \texttt{invalidado} (bool) & \texttt{status}: \texttt{ACTIVE} / \texttt{USED} / \texttt{CANCELLED} \\
Compra fiat & No aplica & Tablas \texttt{orders}, \texttt{order\_items}, \texttt{payments} \\
Compra on-chain & Transferencia atómica SAC XLM + burn/remint & Sincronizado por Indexador vía evento \texttt{boleto\_revendido} \\
Comisiones & Split automático en \texttt{comprar\_boleto} & No almacenado (derivable desde precio y porcentajes del contrato) \\
\bottomrule
\end{tabularx}
\end{table}

\textbf{Nota:} La integración Soroban está completamente implementada. Los campos \texttt{contract\_address}, \texttt{ticket\_root\_id}, \texttt{version}, \texttt{is\_for\_sale}, \texttt{resale\_price} y \texttt{owner\_wallet} en la tabla \texttt{tickets} reflejan el estado on-chain, sincronizado por el Indexador Soroban. La columna \texttt{wallet\_address} en \texttt{users} vincula la identidad Web2 con la dirección Stellar del usuario.

%===============================================================
\FloatBarrier
\clearpage
\section{Atributos de calidad}
\label{sec:calidad}
%===============================================================

Los siguientes escenarios de calidad guían las decisiones arquitectónicas del sistema.

\subsection{QA-01: Seguridad — Autorización por rol}

\begin{table}[htbp]
\centering
\begin{tabularx}{\textwidth}{l X}
\toprule
\textbf{Aspecto} & \textbf{Descripción} \\
\midrule
Estímulo & Un usuario intenta ejecutar una operación para la cual no tiene permisos. \\
Fuente & Actor externo (usuario, contrato, llamada API). \\
Artefacto & Contratos Soroban + Backend Express. \\
Respuesta & La operación es rechazada sin modificar el estado. \\
Medida & 100\% de las funciones que mutan estado validan autorización. \\
\bottomrule
\end{tabularx}
\end{table}

\textbf{Táctica:} On-chain, cada función pública usa \texttt{require\_auth()} sobre la dirección correspondiente (organizador, propietario, verificador, admin). Off-chain, Express usa \texttt{authMiddleware} (valida Bearer JWT y adjunta \texttt{userId}) y los endpoints admin re-consultan el rol desde PostgreSQL en cada solicitud.

\subsection{QA-02: Integridad — Atomicidad en reventa}

\begin{table}[htbp]
\centering
\begin{tabularx}{\textwidth}{l X}
\toprule
\textbf{Aspecto} & \textbf{Descripción} \\
\midrule
Estímulo & Un comprador invoca \texttt{comprar\_boleto} sobre un boleto en reventa. \\
Artefacto & \texttt{ContratoEvento}: función \texttt{comprar\_boleto}. \\
Respuesta & Se transfieren tokens (vendedor + comisiones) y se crea nueva versión del boleto en una sola transacción atómica. Si cualquier paso falla, toda la transacción se revierte. \\
Medida & Cero estados intermedios observables: o se completa todo, o no se modifica nada. \\
\bottomrule
\end{tabularx}
\end{table}

\textbf{Táctica:} Soroban garantiza atomicidad transaccional nativa. El patrón burn/remint invalida la versión anterior y crea la nueva en la misma invocación, eliminando la posibilidad de doble gasto.

\subsection{QA-03: Trazabilidad — Historial de versiones}

\begin{table}[htbp]
\centering
\begin{tabularx}{\textwidth}{l X}
\toprule
\textbf{Aspecto} & \textbf{Descripción} \\
\midrule
Estímulo & Un auditor desea rastrear la cadena de propiedad de un boleto. \\
Artefacto & Storage on-chain del \texttt{ContratoEvento}. \\
Respuesta & Se puede reconstruir el historial completo consultando todas las versiones desde \texttt{(ticket\_root\_id, 0)} hasta \texttt{VersionActual}. \\
Medida & Cada transferencia de propiedad genera una versión inmutable con evento on-chain (\texttt{BoletoRevendido}). \\
\bottomrule
\end{tabularx}
\end{table}

\textbf{Táctica:} El \texttt{ticket\_root\_id} se mantiene constante a través de todas las versiones. Las versiones anteriores permanecen en el storage como registros históricos inmutables. Los eventos on-chain proporcionan un log de auditoría adicional.

\subsection{QA-04: Modificabilidad — Separación on-chain/off-chain}

\begin{table}[htbp]
\centering
\begin{tabularx}{\textwidth}{l X}
\toprule
\textbf{Aspecto} & \textbf{Descripción} \\
\midrule
Estímulo & Se requiere cambiar la interfaz de usuario o agregar un nuevo módulo al backend. \\
Artefacto & Frontend React, Backend Express. \\
Respuesta & El cambio se realiza sin modificar los contratos inteligentes. \\
Medida & Los contratos definen reglas de negocio invariantes; la lógica de presentación y orquestación vive off-chain y es modificable independientemente. \\
\bottomrule
\end{tabularx}
\end{table}

\textbf{Táctica:} Arquitectura en capas con separación estricta. El backend actúa como intermediario y puede evolucionar sin requerir un redeploy de los contratos. El patrón Factory permite desplegar nuevas versiones del \texttt{ContratoEvento} sin afectar los contratos existentes.

\subsection{QA-05: Disponibilidad — Resiliencia del backend}

\begin{table}[htbp]
\centering
\begin{tabularx}{\textwidth}{l X}
\toprule
\textbf{Aspecto} & \textbf{Descripción} \\
\midrule
Estímulo & El servidor backend se reinicia o pierde conexión temporalmente. \\
Artefacto & Backend Express + PostgreSQL (Supabase). \\
Respuesta & Los datos transaccionales persisten en PostgreSQL. Los contratos on-chain siguen operativos independientemente. Al recuperarse el backend, retoma operación normal. \\
Medida & Pérdida de datos cero para transacciones confirmadas en PostgreSQL. Los contratos on-chain tienen disponibilidad inherente de la red Stellar. \\
\bottomrule
\end{tabularx}
\end{table}

\textbf{Táctica:} PostgreSQL provee durabilidad ACID. La blockchain Stellar tiene alta disponibilidad intrínseca (consenso federado, sin downtime planificado). El backend es stateless (sesiones en JWT) y puede reiniciarse sin pérdida de contexto.

%===============================================================
\FloatBarrier
\clearpage
\section{Decisiones arquitectónicas}
\label{sec:decisiones}
%===============================================================

\subsection{ADR-01: Selección de Stellar/Soroban como plataforma blockchain}

\textbf{Contexto:} Se necesita una plataforma blockchain que soporte contratos inteligentes con costos de transacción bajos, finalidad rápida y un ecosistema maduro de tokens.

\textbf{Decisión:} Se selecciona Stellar con Soroban (contratos inteligentes en Rust compilados a WASM) por: transacciones con finalidad en ~5 segundos, costos de gas del orden de fracciones de centavo, soporte nativo de tokens (Stellar Asset Contract), y atomicidad transaccional garantizada.

\textbf{Consecuencias:} El equipo debe aprender Rust y el SDK de Soroban. El ecosistema de herramientas es más joven que el de Ethereum/Solidity, pero suficiente para el alcance del proyecto. Las pruebas se ejecutan con el entorno simulado de Soroban SDK.

\subsection{ADR-02: Patrón Factory para contratos de evento}

\textbf{Contexto:} El sistema debe soportar múltiples eventos, cada uno con su propia configuración de comisiones, token de pago y organizador.

\textbf{Decisión:} Se implementa un contrato \texttt{FabricaBoletos} que despliega instancias independientes de \texttt{ContratoEvento} mediante \texttt{deploy\_v2}. Cada evento tiene su propio contrato con storage aislado.

\textbf{Consecuencias:} Aislamiento total entre eventos (un bug en un contrato no afecta a otros). El storage de Soroban tiene límites por contrato, y este patrón evita cuellos de botella. La factory mantiene un registro de todos los contratos desplegados para descubrimiento. El costo es un deploy adicional por evento.

\subsection{ADR-03: Patrón Burn/Remint para reventa}

\textbf{Contexto:} Al revender un boleto, se necesita transferir propiedad manteniendo trazabilidad completa del historial.

\textbf{Decisión:} En lugar de mutar el propietario del boleto existente, se invalida la versión actual (\texttt{invalidado = true}) y se crea una nueva versión con \texttt{version + 1}, nuevo propietario y nuevo precio. El \texttt{ticket\_root\_id} se mantiene constante.

\textbf{Consecuencias:} Trazabilidad completa: cada versión anterior permanece como registro histórico inmutable en el storage. Se puede reconstruir toda la cadena de propiedad. El costo es mayor uso de storage (una entrada por versión), pero el beneficio en auditoría y prevención de fraude justifica el trade-off.

\subsection{ADR-04: Separación on-chain/off-chain}

\textbf{Contexto:} No toda la lógica del sistema necesita vivir en la blockchain. El almacenamiento on-chain es costoso y las operaciones son lentas comparadas con una base de datos tradicional.

\textbf{Decisión:} Se mantiene on-chain únicamente la lógica crítica de negocio (propiedad, transferencias, comisiones, redención). La gestión de usuarios, catálogo de eventos, carrito de compras, y lógica de presentación vive off-chain en Express + PostgreSQL.

\textbf{Consecuencias:} Mejor experiencia de usuario (operaciones off-chain son instantáneas). Menor costo de gas. La consistencia entre ambos mundos se mantiene mediante el Indexador Soroban, que hace polling cada 5 segundos y sincroniza 7 tipos de eventos on-chain hacia PostgreSQL con auto-recovery de rango de ledgers.

\subsection{ADR-05: Stack tecnológico web}

\textbf{Contexto:} Se necesita un frontend y backend web para la aplicación cliente TuTicket.

\textbf{Decisión:} Frontend con React 18 + Vite + TypeScript + Tailwind CSS + shadcn/ui + \texttt{@stellar/freighter-api}. Backend con Express 4 + TypeScript + Prisma ORM + PostgreSQL (Supabase) + \texttt{@stellar/stellar-sdk}. Autenticación con JWT de token único (7 días de expiración, bcrypt para contraseñas).

\textbf{Consecuencias:} TypeScript compartido entre frontend y backend reduce el cambio de contexto. Prisma genera tipos automáticamente a partir del schema. shadcn/ui acelera el desarrollo de UI con componentes accesibles. Express permite un servidor monolítico simple con el Indexador Soroban corriendo en el mismo proceso, adecuado para la escala del prototipo de tesis.

\subsection{ADR-06: Verificadores autorizados para redención}

\textbf{Contexto:} Se necesita un mecanismo para redimir boletos en la puerta del evento que sea seguro y no dependa de que el propietario auto-redima su boleto (lo cual podría ser manipulado).

\textbf{Decisión:} El organizador registra direcciones de verificadores autorizados on-chain (\texttt{agregar\_verificador}). Solo estos verificadores pueden invocar \texttt{redimir\_boleto}. El propietario del boleto no puede auto-redimirse.

\textbf{Consecuencias:} Separación de responsabilidades clara: el propietario posee el boleto, pero no controla su redención. Previene fraudes donde un propietario podría intentar redimir y luego reclamar que no lo hizo. Los verificadores son gestionables (agregar/remover) por el organizador.

%===============================================================
% Referencias
%===============================================================

\section*{Referencias}
\label{sec:referencias}

\begin{enumerate}[leftmargin=*,label={[\arabic*]}]
  \item IEEE Std 1471-2000: \textit{IEEE Recommended Practice for Architectural Description of Software-Intensive Systems}.
  \item ISO/IEC/IEEE 42010:2011: \textit{Systems and software engineering — Architecture description}.
  \item Brown, S. (2018). \textit{The C4 Model for Visualising Software Architecture}. \url{https://c4model.com/}.
  \item Bass, L., Clements, P., \& Kazman, R. (2012). \textit{Software Architecture in Practice}, 3rd Edition. Addison-Wesley.
  \item SRS del sistema Tickets-NFT B2B (v0.1).
  \item SPMP del sistema Tickets-NFT B2B (v1.1).
  \item Stellar Development Foundation: \textit{Stellar Documentation}, \url{https://developers.stellar.org/}.
  \item Soroban Documentation: \textit{Smart Contracts on Stellar}, \url{https://soroban.stellar.org/}.
\end{enumerate}

\end{document}
